% Source: https://artofproblemsolving.com/community/c3h1765635_counting_triangles
% Stable identity: aops-topic-1765635-q01
Three segments $a<b<c$ form a triangle exactly when $a+b>c$ (the triangle inequality for the longest side). Each length here is the sum of the two before it, so for any choice $a<b<c$ from the list,
\[
a+b\le(\text{the length just before }b)+b=(\text{the length just after }b)\le c .
\]
For example $5+8=13$, $3+5=8$, $3+8<13$. So no choice works:

\begin{asy}
usepackage("amsmath");
size(6cm);
defaultpen(fontsize(10pt));
pair A=(0,0),B=(13,0);
draw(A--B,linewidth(1));
draw(arc(A,5,0,80),blue+dashed);draw(arc(B,8,100,180),red+dashed);
draw(A--5*dir(60),blue);draw(B--B+8*dir(150),red);
dot(A);dot(B);
label("$5$",2.5*dir(60),NW,blue);label("$8$",B+4*dir(150),NE,red);
label("$13$",(6.5,0),S);
label("\mbox{}",(min(currentpicture,user=true)+max(currentpicture,user=true))/2,black); // keeps the text after the figure black when LaTeX typesets the labels inline
\end{asy}

the segments of lengths $5$ and $8$ swung from the ends of the segment of length $13$ can only meet on it. The number of triangles is $\boxed{0}$.
